\documentclass{article}
\usepackage[UTF8]{ctex}
\usepackage{amsmath,amssymb,amsfonts}
\usepackage{enumitem}
\usepackage{geometry}
\geometry{a4paper, margin=1in}
% === 错题本排版优化（不参与题目标识） ===
\setlength{\parindent}{0pt}
\setlength{\parskip}{0.35em plus 0.10em minus 0.08em}
\setlength{\emergencystretch}{3em}
\setlist{itemsep=0.25em, topsep=0.40em, parsep=0pt}
\allowdisplaybreaks[2]
% === 错题本排版优化结束 ===
\begin{document}
\section*{西南财经大学2025年数学分析真题}
\section*{174．西南财经大学2025年数学分析试卷}

\subsection{题目1}
1．求极限 \(\lim _{x \rightarrow+\infty} e^{-x}\left(1+\frac{1}{x}\right)^{x^{2}}\) ．

\subsection{题目2}
2．判断下面命题是否正确，并证明你的结论。\\
函数 \(f(x)\) 在 \((a, b)\) 内处处有导数 \(f^{\prime}(x)\) ，则 \((a, b)\) 中的点不可能为 \(f^{\prime}(x)\) 的第一类间断点．

\begin{itemize}

\subsection{题目3}
\end{itemize}

\subsection{题目3}
原题边界“3．”：函数 \(f(x)\) 在 \([a, b](a b>0)\) 上连续，在 \((a, b)\) 上可微．证明：存在 \(c \in(a, b)\) ，使得
\[
\frac{1}{b-a}\left|\begin{array}{cc}
b & a \\
f(b) & f(a)
\end{array}\right|=f(c)-c f^{\prime}(c) .
\]

\subsection{题目4}

\subsection{题目4}
原题边界“4．”：判断积分 \(\int_{1}^{+\infty} \frac{\sin ^{4} x}{x^{p}\left(x^{p}+\sin x\right)} \mathrm{d} x(p>0)\) 的敛散性。

\subsection{题目5}
5．判断级数 \(\sum_{n=2}^{\infty} \frac{n \cos (n \pi)}{\sqrt{n^{3}-2 n+1}}\) 是否收敛，并指出是绝对收敛还是条件收敛。

\subsection{题目6}
6．已知 \(f(x, y)=\frac{1}{x y}, r=\sqrt{x^{2}+y^{2}}, k>1\) ．

\[
D_{1}=\left\{(x, y) \left\lvert\, \frac{1}{k} x<y<k x\right.\right\}, D_{2}=\{(x, y) \mid x>0, y>0\} .
\]

当 \(i=1,2\) 时，分别判断 \(\lim _{\substack{r \rightarrow+\infty \\(x, y) \in D_{i}}} f(x, y)\) 是否存在．

\subsection{题目7}
7．计算二重积分

\[
\iint_{x^{2}+y^{2} \leq \frac{3}{16}} \min \left\{\sqrt{\frac{3}{16}-x^{2}-y^{2}}, 2\left(x^{2}+y^{2}\right)\right\} \mathrm{d} x \mathrm{~d} y .
\]

\subsection{题目8}
8．已知

\[
f(x, y)= \begin{cases}\frac{x^{m} y^{m}}{\left(x^{2}+y^{2}\right)^{n}} \arcsin \left(x^{2}+y^{2}\right), & (x, y) \neq(0,0) ; \\ 0, & (x, y)=(0,0) .\end{cases}
\]

判断 \(f(x, y)\) 在 \((0,0)\) 处的偏导数是否存在，是否可微．
\end{document}
